\section{Introduction}
Reconstructing the internal properties of transparent and refractive media is a central problem in computer graphics, with applications spanning optical engineering, medical imaging, and visual effects. While opaque surfaces have benefited from the rapid progress of neural rendering techniques, refractive volumes present unique challenges due to non-linear light transport. Standard techniques often approximate light path propagation using straight rays or discrete voxel grids, failing to capture the complex, continuous index of refraction gradients characteristic of real-world materials.

In this work, we propose Vertex-Net, a novel inverse rendering framework designed specifically for complex refractive media. Our method leverages a coordinate-based neural field to parameterize the continuous spatial distribution of the index of refraction. Unlike previous representations that partition space into uniform grids, our continuous neural representation dynamically allocates capacity to high-frequency interfaces. By combining a multi-resolution hash grid with a shallow multi-layer perceptron (MLP), we achieve a highly efficient representation that can be queried at arbitrary spatial coordinates.

To enable optimization from standard multi-view photographs, we integrate this neural field with a differentiable ray tracer that models ray-bending behavior. As light travels through an inhomogeneous medium, its path bends towards regions of higher refractive index, a phenomenon governed by the eikonal equation. Our renderer solves this differential equation numerically, allowing us to compute the derivatives of the rendered pixels with respect to the continuous neural parameters. This differentiable formulation enables end-to-end training of the refractive field using a photometric loss.

In summary, our contributions are: (1) we introduce Vertex-Net, a coordinate-based neural field representation for continuous refractive index reconstruction; (2) we develop a highly optimized, GPU-accelerated differentiable path tracer that handles non-linear ray refraction; and (3) we show that our approach outperforms existing grid-based and volumetric representations on a variety of complex refractive shapes in the Nimbus Refractive Benchmark.
